% Wersja źródłowa: audyt II, jednostka 22
\setcounter{section}{6}
\section{Dystrybucje temperowane i transformacja Fouriera}
\label{sec:temperowane22}

\providecommand{\atlasref}[2]{#2}


\noindent\textbf{Fakty z wcześniejszych rozdziałów.}
Przyjmujemy, że przestrzeń Schwartza jest zupełną przestrzenią Frécheta,
$\mathcal D(\R^d)$ jest w niej gęsta, a transformacja Fouriera
jest jej ciągłym automorfizmem z odwrotnością i wzorem
$\mathcal F^2\varphi=\varphi(-\cdot)$.
Dystrybucja o zwartym nośniku $K$ działa również na funkcjach
$C^\infty$ poprzez $\langle T,f\rangle=\langle T,\chi f\rangle$,
gdzie $\chi=1$ w otoczeniu $K$, i spełnia oszacowanie
skończonego rzędu na ustalonym zwartym otoczeniu $K$.
Korzystamy z Twierdzeń~\ref{thm:schwartz-frechet17},
\ref{thm:density-schwartz17}, \ref{thm:fourier-schwartz17}
i Lematu~\ref{lem:compact-action21}.

W całym rozdziale $d\ge1$ jest liczbą całkowitą.
Rozważamy funkcje oraz dystrybucje zespolone na $\R^d$.
Przestrzeń Schwartza $\mathcal S=\mathcal S(\R^d;\C)$
wyposażamy w topologię Frécheta z rozdziału~17.
Dla $N,m\in\N_0$ przyjmujemy seminormy
$Q_{N,m}(\varphi)=\max_{|\alpha|\le m}\sup_{x\in\R^d}
(1+|x|)^N|\partial^\alpha\varphi(x)|$.
Przez $\mathcal S'$ rozumiemy przestrzeń \emph{ciągłych,
$\C$-liniowych} funkcjonałów na $\mathcal S$.
W parowaniu $\langle T,\varphi\rangle$ nie sprzęgamy zespolenie
funkcji testowej. Przypominamy, że
$\mathcal D(\R^d)=C_c^\infty(\R^d)$ jest gęsta w $\mathcal S$.

\subsection{Dystrybucje temperowane jako przestrzeń dualna}

\noindent\emph{Powiązania tematyczne w Bibliotece matematycznej:} \atlasref{PDE-08::def:pde8-tempered}{PDE-08: definicja dystrybucji temperowanej}.


\begin{definicja}[Dystrybucja temperowana]
\label{def:tempered22}
Dystrybucją temperowaną na $\R^d$ nazywamy element
$T\in\mathcal S'$, czyli ciągły liniowy funkcjonał
$T:\mathcal S\to\C$. Rozpatrujemy na $\mathcal S'$ topologię
słabą-* $\sigma(\mathcal S',\mathcal S)$, wyznaczoną przez
seminormy $T\mapsto|\langle T,\varphi\rangle|$
($\varphi\in\mathcal S$), oraz topologię silną
$\beta(\mathcal S',\mathcal S)$, wyznaczoną przez
$P_B(T)=\sup_{\varphi\in B}|\langle T,\varphi\rangle|$,
gdzie $B$ przebiega wszystkie zbiory ograniczone w $\mathcal S$.
\end{definicja}

\begin{twierdzenie}[Kryterium wzrostu seminormowego]
\label{thm:tempered-estimate22}
Niech $T:\mathcal S(\R^d;\C)\to\C$ będzie liniowy.
Następujące warunki są równoważne:
\begin{enumerate}
\item $T\in\mathcal S'(\R^d)$;
\item istnieją $C>0$ oraz $N,m\in\N_0$ takie, że
$|\langle T,\varphi\rangle|\le C Q_{N,m}(\varphi)$
dla wszystkich $\varphi\in\mathcal S$.
\end{enumerate}
W szczególności ograniczenie $T|_{\mathcal D}$ jest dystrybucją,
a odwzorowanie $\mathcal S'\to\mathcal D'$ dane przez
ograniczenie jest liniowe i różnowartościowe.
\end{twierdzenie}
\begin{proof}
Rodzina $Q_{j,j}$, $j\in\N_0$, jest rosnąca
i wyznacza topologię $\mathcal S$. Jeśli $T$ jest ciągły,
to dla pewnych $j$ i $\varepsilon>0$ warunek
$Q_{j,j}(\varphi)<\varepsilon$ implikuje
$|T(\varphi)|<1$. Liniowość i skalowanie dają
$|T(\varphi)|\le (2/\varepsilon)Q_{j,j}(\varphi)$;
gdy seminorma znika, wniosek wynika ze skalowania przez
dowolną liczbę. Odwrotna implikacja jest natychmiastowa.
Włożenie $j:\mathcal D\hookrightarrow\mathcal S$ jest ciągłe:
dla niepustego zwartego $K\subseteq\R^d$ i
$\varphi\in\mathcal D_K$ mamy
$Q_{N,m}(\varphi)\le (1+\max_{x\in K}|x|)^N p_m(\varphi)$.
Dla $K=\varnothing$ etap jest przestrzenią zerową.
Ciągłość $j$ wynika z własności uniwersalnej granicy
induktywnej (Twierdzenie~\ref{thm:LF-universal18}).
Stąd ograniczenie $T|_{\mathcal D}=T\circ j$ jest ciągłe.
Odwzorowanie $j^{\mathrm t}:T\mapsto T\circ j$ jest
liniowe i różnowartościowe: jeśli dwa ciągłe funkcjonały
zgadzają się na gęstej podprzestrzeni $\mathcal D$,
to ich różnica znika na całej $\mathcal S$.
\end{proof}

\begin{przyklad}[Dystrybucje regularne i miary]
\label{ex:tempered-regular22}
Jeżeli $f\in L^1_{\mathrm{loc}}(\R^d)$ i
dla pewnych $C>0$, $M\ge0$ zachodzi
$|f(x)|\le C(1+|x|)^M$ prawie wszędzie,
to $T_f(\varphi)=\int_{\R^d}f(x)\varphi(x)\,\mathrm d\lambda_d(x)$
określa element $\mathcal S'$. Istotnie, dla całkowitego $N>M+d$ mamy
\[
 \int_{\R^d}|f(x)\varphi(x)|\,\mathrm d\lambda_d(x)
 \le C Q_{N,0}(\varphi)
 \int_{\R^d}(1+|x|)^{M-N}\,\mathrm d\lambda_d(x).
\]
Ostatnia całka jest skończona: na kuli jednostkowej
funkcja jest ograniczona, a na pierścieniu
$2^k\le |x|<2^{k+1}$ całkę ogranicza stała razy
$2^{k(d+M-N)}$. Szereg tych oszacowań jest zbieżny.
Z Twierdzenia~\ref{thm:tempered-estimate22} wynika
temperowalność $T_f$.
Bardziej ogólnie, jeżeli $\mu$ jest zespoloną lokalnie skończoną
miarą borelowską i dla pewnego $N\ge0$ zachodzi
$\int_{\R^d}(1+|x|)^{-N}\,\mathrm d|\mu|(x)<\infty$,
to $\langle T_\mu,\varphi\rangle=
\int_{\R^d}\varphi(x)\,\mathrm d\mu(x)$
definiuje dystrybucję temperowaną. Istotnie, biorąc
liczbę całkowitą $n\ge N$, mamy
\[
 \left|\int_{\R^d}\varphi\,\mathrm d\mu\right|
 \le Q_{n,0}(\varphi)
 \int_{\R^d}(1+|x|)^{-N}\,\mathrm d|\mu|(x).
\]
Całka jest bezwzględnie zbieżna względem $|\mu|$,
a oszacowanie dowodzi ciągłości funkcjonału.
Przykładami są funkcja stała $1$, wielomiany i miary Diraca.
\end{przyklad}

\begin{przyklad}[Dystrybucja nietemperowana]
\label{ex:non-tempered22}
Dystrybucja regularna wyznaczona przez
$f(x)=e^{|x|^2}$ należy do $\mathcal D'(\R^d)$,
ale nie jest ograniczeniem żadnego elementu $\mathcal S'$.
Aby to wykazać, weźmy $\chi\in\mathcal D(\R^d)$,
$0\le\chi\le1$, równą $1$ na $\overline B(0,1)$,
i funkcje $\varphi_R(x)=\chi(x/R)e^{-|x|^2/2}$.
Dla $R\to\infty$ zachodzi
$\varphi_R\to e^{-|x|^2/2}$ w $\mathcal S$
na mocy dowodu Twierdzenia~\ref{thm:density-schwartz17}:
dla $g(x)=e^{-|x|^2/2}\in\mathcal S$ i $R\ge1$ zachodzi
$Q_{N,m}(g-\chi(\cdot/R)g)\le
C_m(1+R)^{-1}Q_{N+1,m}(g)$.
Przynależność $g$ do $\mathcal S$ wynika stąd, że
każda jego pochodna jest wielomianem pomnożonym przez $g$.
Jednak
$\langle T_f,\varphi_R\rangle
=\int_{\R^d}\chi(x/R)e^{|x|^2/2}\,\mathrm d\lambda_d(x)
\to+\infty$,
gdyż dla $R\ge1$ całka jest nie mniejsza niż
$\int_{B(0,R)}1\,\mathrm d\lambda_d=
R^d\lambda_d(B(0,1))$.
Gdyby istniało ciągłe przedłużenie na $\mathcal S$,
wartości te musiałyby zbiegać do skończonej liczby.
\end{przyklad}

\begin{stwierdzenie}[Podprzestrzenie klasycznych dystrybucji]
\label{prop:schwartz-embeddings22}
Istnieją naturalne, różnowartościowe i ciągłe odwzorowania
\[
 \mathcal D\hookrightarrow\mathcal S\hookrightarrow
 L^2(\R^d)\hookrightarrow\mathcal S'\hookrightarrow\mathcal D',
\]
przy czym obie przestrzenie dualne wyposażamy albo
w topologie silne, albo w topologie słabe-*,
gdzie funkcję $f\in L^2$ identyfikujemy z funkcjonałem
$\langle T_f,\varphi\rangle=\int_{\R^d}f(x)\varphi(x)
\,\mathrm d\lambda_d(x)$.
Ponadto każda dystrybucja o zwartym nośniku jest temperowana:
$\mathcal E'(\R^d)\subseteq\mathcal S'(\R^d)$.
\end{stwierdzenie}
\begin{proof}
Pierwsze włożenie wykazano w Twierdzeniu~
\ref{thm:tempered-estimate22}. Drugie jest ciągłe na mocy
Stwierdzenia~\ref{prop:schwartz-operators17} i różnowartościowe,
ponieważ ciągła funkcja równa zeru prawie wszędzie znika
wszędzie. Dla $f\in L^2$ nierówność Cauchy'ego--Schwarza daje
$|\langle T_f,\varphi\rangle|\le\|f\|_2\|\varphi\|_2$;
także $\|\varphi\|_2\le C_dQ_{d+1,0}(\varphi)$.
Dla ograniczonego $B\subseteq\mathcal S$ otrzymujemy
\[
 P_B(T_f)\le C_d\|f\|_2\sup_{\varphi\in B}Q_{d+1,0}(\varphi).
\]
Prawa strona jest skończona i dowodzi ciągłości silnej,
a dla singletonów także słabej-*.
Jeśli $T_f=0$, to $\int f\varphi\,\mathrm d\lambda_d=0$
dla wszystkich $\varphi\in\mathcal D$.
Z gęstości $\mathcal D$ w $L^2$ wybieramy $\varphi_n\to
\overline f$ w $L^2$; nierówność Cauchy'ego--Schwarza daje
$\int |f|^2\,\mathrm d\lambda_d=0$, więc $f=0$ w $L^2$.
Ostatnie odwzorowanie jest różnowartościowe z
Twierdzenia~\ref{thm:tempered-estimate22}. Dla
$B\subseteq\mathcal D$ ograniczonego mamy
$P_B(T|_{\mathcal D})=P_{j(B)}(T)$, gdzie $j(B)$ jest
ograniczony w $\mathcal S$ dzięki ciągłości $j$.
Daje to ciągłość silną, a oceny na pojedynczych
funkcjach testowych dają ciągłość słabą-*.
Jeżeli $T\in\mathcal E'$ ma zwarty nośnik,
to jego przedłużenie z Lematu~\ref{lem:compact-action21}
spełnia $|T(\varphi)|\le C Q_{0,m}(\varphi)$.
Jest więc ciągłe na $\mathcal S$ i zgadza się z $T$
na $\mathcal D$; identyfikujemy $T$ z tym jednoznacznym
przedłużeniem.
\end{proof}

\subsection{Operatory na przestrzeni dystrybucji temperowanych}

\begin{twierdzenie}[Różniczkowanie i mnożenie przez wielomian]
\label{thm:tempered-operators22}
Dla każdego $\alpha,\beta\in\N_0^d$ operatory
$T\mapsto\partial^\alpha T$ i $T\mapsto x^\beta T$
są liniowe i ciągłe na $\mathcal S'$ zarówno dla topologii
$\sigma(\mathcal S',\mathcal S)$, jak i
$\beta(\mathcal S',\mathcal S)$, gdzie
\[
\langle\partial^\alpha T,\varphi\rangle
=(-1)^{|\alpha|}\langle T,\partial^\alpha\varphi\rangle,
\qquad\langle x^\beta T,\varphi\rangle
=\langle T,x^\beta\varphi\rangle.
\]
To samo dotyczy mnożenia przez każdą funkcję
$a\in C^\infty(\R^d;\C)$, dla której dla każdego
$\gamma\in\N_0^d$ istnieją $C_\gamma>0$ oraz
$M_\gamma\in\N_0$ takie, że
$|\partial^\gamma a(x)|\le C_\gamma(1+|x|)^{M_\gamma}$
dla wszystkich $x\in\R^d$.
\end{twierdzenie}
\begin{proof}
Różniczkowanie i mnożenie przez $x^\beta$ są ciągłymi
operatorami liniowymi $\mathcal S\to\mathcal S$,
co wykazano w Stwierdzeniu~\ref{prop:schwartz-operators17}.
Dla ogólnego mnożnika $a$ ustalmy $N,m\in\N_0$,
$M_m=\max_{|\gamma|\le m}M_\gamma$ i
$C_m=\max_{|\beta|\le m}\sum_{\gamma\le\beta}
\binom\beta\gamma C_\gamma$.
Ze wzoru Leibniza, dla $|\beta|\le m$, dostajemy
\[
 (1+|x|)^N|\partial^\beta(a\varphi)(x)|
 \le C_m Q_{N+M_m,m}(\varphi).
\]
Po wzięciu supremum oraz maksimum otrzymujemy
$Q_{N,m}(a\varphi)\le C_m Q_{N+M_m,m}(\varphi)$.
Zatem odwzorowanie $\varphi\mapsto a\varphi$ jest
ciągłym operatorem na $\mathcal S$.
Dla dowolnego ciągłego operatora liniowego
$A:\mathcal S\to\mathcal S$ jego transpozycja
$A^{\mathrm t}T=T\circ A$ spełnia
$P_B(A^{\mathrm t}T)=P_{A(B)}(T)$.
Ciągłość $A$ zapewnia ograniczoność $A(B)$, skąd
ciągłość silna; dla singletonów otrzymujemy ciągłość
słabą-*. Stosujemy to do $A\varphi=(-1)^{|\alpha|}
\partial^\alpha\varphi$ i $A\varphi=a\varphi$,
zgodnie z definicjami pochodnej i mnożenia dystrybucji.
Ograniczenie tych działań do $\mathcal D$ zgadza się
z działaniami na zwykłych dystrybucjach.
\end{proof}

\begin{lemat}[Względna zwartość zbiorów ograniczonych w $\mathcal S$]
\label{lem:montel-schwartz22}
Każdy zbiór ograniczony w $\mathcal S(\R^d)$ ma zwarte
domknięcie w jej topologii Frécheta.
\end{lemat}
\begin{proof}
Niech $B\subseteq\mathcal S$ będzie ograniczony i
$(\varphi_n)\subseteq B$ dowolnym ciągiem.
Dla każdej liczby całkowitej $R\ge1$ i każdego $m\ge0$
rodziny pochodnych $\partial^\alpha\varphi_n$ dla
$|\alpha|\le m$ są jednostajnie ograniczone i
jednostajnie równociągłe na $\overline B(0,R)$:
równociągłość wynika z ograniczenia pochodnych
rzędu $m+1$ i całkowania wzdłuż odcinka w kuli:
dla $x,y$ w tej kuli oraz $|\alpha|\le m$ mamy
$|\partial^\alpha\varphi_n(x)-\partial^\alpha\varphi_n(y)|
\le\|x-y\|_1\sup_{\psi\in B}Q_{0,m+1}(\psi)$. Z twierdzenia Arzelà--Ascoliego
i argumentu diagonalnego uzyskujemy podciąg
$(\varphi_{n_k})$, którego wszystkie pochodne zbiegają
jednostajnie na wszystkich kulach domkniętych.
Nie musimy osobno identyfikować granic pochodnych:
poniższe oszacowanie da zbieżność w zupełnej
przestrzeni Schwartza.
Dla ustalonych $N,m$ i $R\ge1$ poza $B(0,R)$ zachodzi
\[
(1+|x|)^N|\partial^\alpha(\varphi_{n_k}-\varphi_{n_l})(x)|
\le R^{-1}\bigl(Q_{N+1,m}(\varphi_{n_k})
+Q_{N+1,m}(\varphi_{n_l})\bigr)\quad(|\alpha|\le m).
\]
Prawa strona jest jednostajnie mała przy dużym $R$;
wewnątrz kuli $B(0,R)$ zbieżność wszystkich pochodnych
jest jednostajna. Dlatego $(\varphi_{n_k})$ jest ciągiem
Cauchy'ego we wszystkich seminormach $Q_{N,m}$.
Zupełność $\mathcal S$ daje zbieżność do elementu
$\mathcal S$. Ponieważ przestrzeń jest metryzowalna,
ciągowa względna zwartość zbioru jest równoważna
zwartości jego domknięcia. Istotnie, dla ciągu $(g_n)$
w domknięciu $B$ wybieramy $f_n\in B$ w odległości
metrycznej mniejszej niż $1/n$ od $g_n$.
Zbieżny podciąg $(f_{n_k})$ daje zbieżny podciąg
$(g_{n_k})$ o tej samej granicy, należącej do domknięcia.
Zwartość wynika z kryterium ciągowego dla przestrzeni
metrycznych.
\end{proof}

\begin{twierdzenie}[Zbieżność słaba-* i silna ciągów]
\label{thm:seq-topologies-schwartz22}
Niech $(T_n)$ będzie ciągiem dystrybucji temperowanych,
a $T\in\mathcal S'$. Wtedy
\[
T_n\longrightarrow T\text{ w }\sigma(\mathcal S',\mathcal S)
\quad\Longleftrightarrow\quad
T_n\longrightarrow T\text{ w }\beta(\mathcal S',\mathcal S).
\]
Mimo to topologie $\sigma(\mathcal S',\mathcal S)$ i
$\beta(\mathcal S',\mathcal S)$ nie są równe.
\end{twierdzenie}
\begin{proof}
Implikacja ze zbieżności silnej jest natychmiastowa.
Załóżmy zbieżność słabą-*. Rodzina
$\{T_n-T:n\ge1\}$ jest punktowo ograniczona.
Ponieważ $\mathcal S$ jest przestrzenią Frécheta,
a więc beczkową, twierdzenie Banacha--Steinhausa
implikuje jej równociągłość.
Dla dowolnego ograniczonego $B\subseteq\mathcal S$
Lemat~\ref{lem:montel-schwartz22} daje zwartość
$K=\overline B$. Równociągłość funkcjonałów i
punktowa zbieżność do zera zapewniają zbieżność
jednostajną na $K$: dla danego $\varepsilon>0$
wybieramy otoczenie $V$ z
$\sup_n|\langle T_n-T,v\rangle|<\varepsilon/2$
dla $v\in V$ i skończone pokrycie $K$ przesunięciami
$\varphi_j+V$. Następnie dla dostatecznie dużego $n$
wszystkie $|\langle T_n-T,\varphi_j\rangle|<\varepsilon/2$.
Wówczas $P_B(T_n-T)<\varepsilon$.

Aby wykazać różność topologii, wybierzmy niezerową funkcję
$\psi\in\mathcal D$ o nośniku w $B(0,1)$ i zdefiniujmy
$\psi_n(x)=e^{-n^2}\psi(x-3ne_1)$ dla $n\ge1$.
Zbiór $B=\{\psi_n:n\ge1\}$ jest ograniczony w
$\mathcal S$: dla ustalonych $N,m$ mamy
$Q_{N,m}(\psi_n)\le e^{-n^2}(2+3n)^N Q_{0,m}(\psi)$,
co jest ograniczone niezależnie od $n$.
Jego liniowa otoczka jest nieskończenie wymiarowa,
ponieważ nośniki są parami rozłączne i każda $\psi_n$
jest niezerowa.
Otoczenie silne $\{T:P_B(T)<1/2\}$ nie zawiera
żadnego słabo-* otoczenia zera: dla dowolnie wielu
$\varphi_1,\ldots,\varphi_r\in\mathcal S$
można wybrać $\psi_n$ poza ich liniową otoczką.
Podprzestrzeń $M=\operatorname{span}\{\varphi_1,\ldots,
\varphi_r\}$ jest skończenie wymiarowa, a więc domknięta
w przestrzeni lokalnie wypukłej Hausdorffa $\mathcal S$.
Zatem $\mathcal S/M$ jest przestrzenią lokalnie wypukłą
Hausdorffa, a klasa $\psi_n+M$ jest niezerowa.
Twierdzenie Hahna--Banacha daje ciągły funkcjonał
na ilorazie niezerowy na tej klasie; po normalizacji
i złożeniu z odwzorowaniem ilorazowym otrzymujemy
$T\in\mathcal S'$ taki, że $T(\varphi_j)=0$ dla
wszystkich $j$ oraz $T(\psi_n)=1$.
Należy on do dowolnego podstawowego słabo-* otoczenia
wyznaczonego przez te $\varphi_j$, ale nie do
$\{T:P_B(T)<1/2\}$. Ponieważ topologia silna jest
co najmniej tak mocna jak słaba-*, dowodzi to ich różności.
\end{proof}

\subsection{Rozszerzenie transformacji Fouriera}

\noindent\emph{Powiązania tematyczne w Bibliotece matematycznej:} \atlasref{PDE-08::def:pde8-distFT}{PDE-08: transformacja Fouriera dystrybucji}.


\begin{definicja}[Transformacja Fouriera dystrybucji temperowanej]
\label{def:fourier-tempered22}
Dla $T\in\mathcal S'(\R^d)$ definiujemy
$\widehat T=\mathcal F T\in\mathcal S'$ przez
\[
\langle\mathcal F T,\varphi\rangle
=\langle T,\mathcal F\varphi\rangle
\qquad(\varphi\in\mathcal S).
\]
Występuje tu transformacja Fouriera na $\mathcal S$
z Definicji~\ref{def:fourier-schwartz17}, z normalizacją
$(2\pi)^{-d/2}$ i jądrem $e^{-ix\cdot\xi}$.
Ponieważ $\mathcal F:\mathcal S\to\mathcal S$ jest ciągła
(Twierdzenie~\ref{thm:fourier-schwartz17}), złożenie
$T\circ\mathcal F$ jest ciągłym $\C$-liniowym funkcjonałem.
Definicja nie zawiera sprzężenia zespolonego ani
transformacji odwrotnej.
\end{definicja}

\begin{twierdzenie}[Transformacja Fouriera jest automorfizmem]
\label{thm:fourier-auto-tempered22}
Transformacja $\mathcal F:\mathcal S'\to\mathcal S'$
jest ciągłym izomorfizmem liniowo-topologicznym
zarówno dla topologii $\sigma(\mathcal S',\mathcal S)$,
jak i $\beta(\mathcal S',\mathcal S)$.
Jej odwrotność dana jest transpozycją
$\mathcal F^{-1}$ na $\mathcal S$:
$\langle\mathcal F^{-1}T,\varphi\rangle
=\langle T,\mathcal F^{-1}\varphi\rangle$.
Dla każdego $T\in\mathcal S'$ i $j\in\{1,\ldots,d\}$
zachodzą wzory
\[
\mathcal F(\partial_jT)=i\xi_j\,\mathcal F T,
\qquad
\mathcal F(x_jT)=i\partial_{\xi_j}(\mathcal F T),
\qquad\mathcal F^2T=T(-\cdot),
\]
gdzie $\langle T(-\cdot),\varphi\rangle=
\langle T,\varphi(-\cdot)\rangle$.
\end{twierdzenie}
\begin{proof}
Z Twierdzenia~\ref{thm:fourier-schwartz17} oba operatory
$\mathcal F$ i $\mathcal F^{-1}$ na $\mathcal S$ są
ciągłe. Dla $\varphi\in\mathcal S$ mamy
$\langle\mathcal F^{-1}\mathcal FT,\varphi\rangle
=\langle T,\mathcal F\mathcal F^{-1}\varphi\rangle
=\langle T,\varphi\rangle$.
Analogicznie
$\langle\mathcal F\mathcal F^{-1}T,\varphi\rangle
=\langle T,\mathcal F^{-1}\mathcal F\varphi\rangle
=\langle T,\varphi\rangle$.
Są to zatem wzajemnie odwrotne operatory na $\mathcal S'$.
Ciągłość słaba-* wynika z działania na pojedynczej
funkcji testowej. Dla silnej mamy
$P_B(\mathcal FT)=P_{\mathcal F(B)}(T)$;
$\mathcal F(B)$ jest ograniczony, gdy $B$ jest ograniczony,
co daje ciągłość. Odwrotność rozpatrujemy analogicznie.
Ustalmy $\varphi\in\mathcal S$, oznaczając jej argument
przez $\xi$, a argument $\mathcal F\varphi$ przez $x$.
Z Twierdzenia~\ref{thm:fourier-schwartz17} wynika
$\partial_{x_j}\mathcal F\varphi
=\mathcal F(-i\xi_j\varphi)$ oraz
$x_j\mathcal F\varphi=\mathcal F(-i\partial_{\xi_j}\varphi)$.
Zatem
\[
\begin{aligned}
\langle\mathcal F(\partial_jT),\varphi\rangle
&=-\langle T,\partial_{x_j}\mathcal F\varphi\rangle\\
&=\langle T,\mathcal F(i\xi_j\varphi)\rangle
=\langle i\xi_j\mathcal FT,\varphi\rangle,\\
\langle\mathcal F(x_jT),\varphi\rangle
&=\langle T,x_j\mathcal F\varphi\rangle\\
&=\langle T,\mathcal F(-i\partial_{\xi_j}\varphi)\rangle\\
&=-i\langle\mathcal FT,\partial_{\xi_j}\varphi\rangle
=\langle i\partial_{\xi_j}\mathcal FT,\varphi\rangle.
\end{aligned}
\]
Wzór odwrotny na $\mathcal S$ daje
$\mathcal F^{-1}\varphi(x)=\mathcal F\varphi(-x)$,
więc $\mathcal F^2\varphi(x)=\varphi(-x)$.
W konsekwencji
$\langle\mathcal F^2T,\varphi\rangle
=\langle T,\mathcal F^2\varphi\rangle
=\langle T,\varphi(-\cdot)\rangle$.
W szczególności $\mathcal F^4T=T$.
\end{proof}

\begin{przyklad}[Funkcja stała, delta Diraca i pochodne]
\label{ex:fourier-delta22}
Dla $a\in\R^d$ i $\alpha\in\N_0^d$ mamy w $\mathcal S'$
wzory
\[
\begin{aligned}
\mathcal F\delta_a&=(2\pi)^{-d/2}T_{e^{-ia\cdot\xi}},\\
\mathcal F T_1&=(2\pi)^{d/2}\delta_0,\\
\mathcal F(\partial^\alpha\delta_0)
&=(2\pi)^{-d/2}T_{(i\xi)^\alpha}.
\end{aligned}
\]
Przez $T_f$ oznaczamy dystrybucję regularną
reprezentowaną przez $f$, w szczególności $T_1$
odpowiada funkcji stałej $1$.
Pierwszy wynika z
$\langle\mathcal F\delta_a,\varphi\rangle
=(2\pi)^{-d/2}\int e^{-ia\cdot\xi}\varphi(\xi)
\,\mathrm d\lambda_d(\xi)$.
Dla $a=0$ mamy $\mathcal F\delta_0=(2\pi)^{-d/2}T_1$.
Ponieważ odbicie $\delta_0$ jest równe $\delta_0$,
zastosowanie $\mathcal F$ daje
$\delta_0=(2\pi)^{-d/2}\mathcal F T_1$.
Iterując wzór na transformatę pochodnej z
Twierdzenia~\ref{thm:fourier-auto-tempered22}, dostajemy
$\mathcal F(\partial^\alpha\delta_0)
=(i\xi)^\alpha\mathcal F\delta_0$, co daje trzeci wzór.
W szczególności delta Diraca nie jest funkcją,
ale jej transformata Fouriera jest funkcją gładką.
\end{przyklad}

\begin{twierdzenie}[Twierdzenie Plancherela]
\label{thm:plancherel22}
Transformacja Fouriera z normalizacją $(2\pi)^{-d/2}$
przedłuża się jednoznacznie do operatora unitarnego
$\mathcal F:L^2(\R^d;\C)\to L^2(\R^d;\C)$,
czyli dla $f,g\in L^2$ zachodzi
$\int_{\R^d}\widehat f(\xi)\overline{\widehat g(\xi)}
\,\mathrm d\lambda_d(\xi)
=\int_{\R^d}f(x)\overline{g(x)}\,\mathrm d\lambda_d(x)$.
Dla $f\in L^2$ zachodzi ponadto
$\mathcal F T_f=T_{\widehat f}$, gdzie $\widehat f$
oznacza transformatę w $L^2$.
\end{twierdzenie}
\begin{proof}
Najpierw ustalmy $f,\varphi\in\mathcal S$.
Ponieważ $f,\varphi\in L^1$, moduł funkcji
podcałkowej w całce podwójnej jest całkowalny:
\[
 \int_{\R^d}\int_{\R^d}|f(x)\varphi(\xi)|
 \,\mathrm d\lambda_d(x)\,\mathrm d\lambda_d(\xi)
 =\|f\|_1\|\varphi\|_1.
\]
Twierdzenie Fubiniego i symetria jądra
$e^{-ix\cdot\xi}$ dają
\[
 \int_{\R^d}\widehat f(\xi)\varphi(\xi)
 \,\mathrm d\lambda_d(\xi)
 =\int_{\R^d}f(x)(\mathcal F\varphi)(x)
 \,\mathrm d\lambda_d(x).
\]
Jest to zgodność $\mathcal F T_f=T_{\widehat f}$ na
$\mathcal S$; parowanie jest biliniowe, a nie
półtoraliniowe.

Niech teraz $f,g\in\mathcal S$. Stosujemy powyższą
tożsamość do $\varphi=\overline{\widehat g}\in\mathcal S$.
Ze wzoru odwrotnego w Twierdzeniu~\ref{thm:fourier-schwartz17}
wynika
$\mathcal F(\overline{\widehat g})(x)
=\overline{(\mathcal F^{-1}\widehat g)(x)}
=\overline{g(x)}$.
Otrzymujemy
\[
 \int_{\R^d}\widehat f(\xi)\overline{\widehat g(\xi)}
 \,\mathrm d\lambda_d(\xi)
 =\int_{\R^d}f(x)\overline{g(x)}\,\mathrm d\lambda_d(x).
\]
W szczególności $\|\widehat f\|_2=\|f\|_2$.
Przestrzeń $\mathcal S$ jest gęsta w $L^2$, bo zawiera
gęstą podprzestrzeń $\mathcal D$. Dla $f\in L^2$
wybieramy $f_n\in\mathcal S$ z $f_n\to f$ w $L^2$.
Ciąg $(\widehat f_n)$ jest Cauchy'ego w $L^2$,
ponieważ $\|\widehat f_n-\widehat f_m\|_2
=\|f_n-f_m\|_2$. Z zupełności istnieje granica;
dla innego ciągu przybliżającego $(g_n)$ tę samą granicę
zapewnia $\|\widehat f_n-\widehat g_n\|_2
=\|f_n-g_n\|_2\to0$. Definiujemy $\widehat f$ jako
tak uzyskaną granicę. Liniowość i zachowanie normy
przechodzą na granicę; jednoznaczność przedłużenia
wynika z gęstości $\mathcal S$.
Obraz izometrii jest domknięty: jeśli ciąg obrazów
jest zbieżny, ciąg przeciwobrazów jest Cauchy'ego i
jego granica ma żądany obraz. Obraz zawiera
$\mathcal F(\mathcal S)=\mathcal S$, więc jest całym
$L^2$. Zachowanie iloczynu skalarnego wynika z ciągłości
iloczynu skalarnego na $L^2$ i tożsamości na $\mathcal S$.
Otrzymujemy zatem operator unitarny.

Na koniec, dla $f_n\to f$ w $L^2$ oraz ograniczonego
$B\subseteq\mathcal S$ mamy
\[
 P_B(T_{f_n}-T_f)
 \le C_d\|f_n-f\|_2
 \sup_{\varphi\in B}Q_{d+1,0}(\varphi)\longrightarrow0.
\]
Analogicznie $T_{\widehat f_n}\to T_{\widehat f}$ silnie
w $\mathcal S'$, bo $\widehat f_n\to\widehat f$ w $L^2$.
Z ciągłości $\mathcal F$ w topologii silnej i
wykazanej zgodności na $\mathcal S$ otrzymujemy
$\mathcal FT_f=\lim_n\mathcal FT_{f_n}
=\lim_nT_{\widehat f_n}=T_{\widehat f}$.
Jednoznaczność granicy wynika z warunku Hausdorffa
dla topologii silnej na $\mathcal S'$.
\end{proof}

\subsection{Splot z funkcją szybko malejącą i transformata dystrybucji o zwartym nośniku}

\begin{twierdzenie}[Splot z elementem przestrzeni Schwartza]
\label{thm:schwartz-convolution22}
Niech $T\in\mathcal S'(\R^d)$ i
$\rho\in\mathcal S(\R^d)$. Funkcja
$u(x)=\langle T_y,\rho(x-y)\rangle$ należy do
$C^\infty(\R^d)$, a każda jej pochodna ma
co najwyżej wzrost wielomianowy. Dystrybucja
regularna $T_u$ jest temperowana i oznaczamy ją $T*\rho$.
Dla każdego $\alpha\in\N_0^d$ zachodzi
\[
\mathcal F(T*\rho)=(2\pi)^{d/2}(\mathcal FT)\,(\mathcal F\rho),
\qquad \partial^\alpha(T*\rho)=T*(\partial^\alpha\rho).
\]
\end{twierdzenie}
\begin{proof}
Dla $h\in\R^d$ oraz $N,m\in\N_0$ zachodzi
$Q_{N,m}(\rho(h-\cdot))\le(1+|h|)^N Q_{N,m}(\rho)$,
gdyż $1+|y|\le(1+|h|)(1+|h-y|)$.
Ta sama nierówność obowiązuje dla każdej pochodnej
$\rho$. Z twierdzenia Newtona--Leibniza zastosowanego
dwukrotnie, dla $0<|t|\le1$ i $j\in\{1,\ldots,d\}$,
otrzymujemy
\[
\begin{aligned}
 Q_{N,m}\left(
 \frac{\rho(x+te_j-\cdot)-\rho(x-\cdot)}{t}
 -(\partial_j\rho)(x-\cdot)\right)
 \le |t|(2+|x|)^N Q_{N,m+2}(\rho).
\end{aligned}
\]
Istotnie, po każdej pochodnej względem $y$ resztę
ograniczamy przez $|t|$ razy supremum odpowiedniej
pochodnej drugiego rzędu względem parametru na
odcinku od $x$ do $x+te_j$ i stosujemy poprzednią
nierówność wagową. Jednokrotne zastosowanie twierdzenia
Newtona--Leibniza daje również ciągłość przesunięć
w każdej seminormie. Powtarzając argument dla
pochodnych $\rho$, otrzymujemy gładkość funkcji
$x\mapsto\rho(x-\cdot)$ o wartościach w $\mathcal S$.
Ciągłość $T$ pozwala więc różniczkować pod parowaniem:
\[
 \partial^\alpha u(x)=\langle T_y,
 (\partial^\alpha\rho)(x-y)\rangle.
\]
Z Twierdzenia~\ref{thm:tempered-estimate22} wybieramy
$C>0$, $N,m\in\N_0$ z $|T(\varphi)|\le C Q_{N,m}(\varphi)$.
Poprzednie oszacowanie przesunięcia daje
\[
 |\partial^\alpha u(x)|\le
 C(1+|x|)^N Q_{N,m+|\alpha|}(\rho).
\]
W szczególności $T_u\in\mathcal S'$ na mocy
Przykładu~\ref{ex:tempered-regular22}.

Ustalmy $\psi\in\mathcal S$ i niech $G_x(y)=
\psi(x)\rho(x-y)$. Odwzorowanie $x\mapsto G_x$
jest ciągłe w $\mathcal S$ i
\[
 Q_{N,m}(G_x)\le |\psi(x)|(1+|x|)^N Q_{N,m}(\rho).
\]
Prawa strona jest całkowalna po $\R^d$ dla każdych $N,m$.
Dla $R\ge1$ określmy $I_R(y)=\int_{[-R,R]^d}G_x(y)
\,\mathrm d\lambda_d(x)$. Różniczkowanie pod całką na
tym zwartym prostokącie jest dozwolone, a powyższe
oszacowanie dowodzi, że $I_R\in\mathcal S$.
Sumy Riemanna tej całki zbiegają do $I_R$ w każdej
seminormie $Q_{N,m}$: błąd jest nie większy niż
objętość prostokąta razy supremum
$Q_{N,m}(G_x-G_{x'})$ dla $x,x'$ z tej samej części
podziału. Supremum dąży do zera z jednostajnej
ciągłości w danej seminormie na prostokącie.
Stosując $T$ do sum Riemanna i przechodząc do granicy
otrzymujemy
\[
 \langle T,I_R\rangle=
 \int_{[-R,R]^d}\psi(x)u(x)\,\mathrm d\lambda_d(x).
\]
Dla $R'<R''$ oszacowanie ogona ma postać
\[
 Q_{N,m}(I_{R''}-I_{R'})
 \le Q_{N,m}(\rho)
 \int_{\R^d\setminus[-R',R']^d}
 (1+|x|)^N|\psi(x)|\,\mathrm d\lambda_d(x).
\]
W szczególności $I_n$ jest ciągiem Cauchy'ego w $\mathcal S$.
Z jej zupełności granica istnieje w $\mathcal S$, a
zbieżność punktowa identyfikuje ją z funkcją
$\check\rho*\psi$, gdzie $\check\rho(y)=\rho(-y)$:
\[
 (\check\rho*\psi)(y)=
 \int_{\R^d}\rho(x-y)\psi(x)\,\mathrm d\lambda_d(x).
\]
Z ciągłości $T$ oraz bezwzględnej całkowalności
$u\psi$ otrzymujemy
$\langle T_u,\psi\rangle=\langle T,\check\rho*\psi\rangle$.
Stwierdzenie~\ref{prop:convolution-schwartz17} zapewnia
ciągłość operatora $\psi\mapsto\check\rho*\psi$ na
$\mathcal S$. Jego transpozycją jest więc ciągły
operator $T\mapsto T*\rho$ dla obu topologii na
$\mathcal S'$, zgodnie z dowodem
Twierdzenia~\ref{thm:tempered-operators22}.
Jeśli $\rho\in\mathcal D$, konstrukcja zgadza się na
$\mathcal D$ z Twierdzeniem~\ref{thm:smoothing21}.

Pozostaje obliczyć stałą i znak w formule Fouriera.
Dla $\psi\in\mathcal S$ mamy, przy $c_d=(2\pi)^{-d/2}$,
\[
\begin{aligned}
(\check\rho*\mathcal F\psi)(y)
 &=\int_{\R^d}\rho(z)(\mathcal F\psi)(y+z)
 \,\mathrm d\lambda_d(z)\\
 &=c_d\int_{\R^d}e^{-iy\cdot\xi}\psi(\xi)
 \int_{\R^d}\rho(z)e^{-iz\cdot\xi}\,\mathrm d\lambda_d(z)
 \,\mathrm d\lambda_d(\xi)\\
 &=(2\pi)^{d/2}\mathcal F((\mathcal F\rho)\psi)(y).
\end{aligned}
\]
Twierdzenie Fubiniego stosujemy, ponieważ całka
modułu podwójnego jest ograniczona przez
$c_d\|\rho\|_1\|\psi\|_1$. Iloczyn
$(\mathcal F\rho)\psi$ należy do $\mathcal S$.
Zatem
\[
\begin{aligned}
\langle\mathcal F(T*\rho),\psi\rangle
&=\langle T,\check\rho*\mathcal F\psi\rangle\\
&=(2\pi)^{d/2}\langle\mathcal FT,
 (\mathcal F\rho)\psi\rangle.
\end{aligned}
\]
To jest pierwszy wzór, a drugi wynika ze wzoru na
$\partial^\alpha u$ i zgodności pochodnej dystrybucji
regularnej z klasyczną pochodną gładkiej funkcji
(wystarczy sprawdzić na $\mathcal D$, gęstej w $\mathcal S$).
\end{proof}

\begin{twierdzenie}[Transformata dystrybucji o zwartym nośniku]
\label{thm:compact-fourier22}
Jeżeli $T\in\mathcal E'(\R^d)$ ma zwarty nośnik $K$,
to jego transformata Fouriera jest dystrybucją regularną
reprezentowaną przez gładką funkcję
\[
F_T(\xi)=(2\pi)^{-d/2}\langle T_x,e^{-ix\cdot\xi}\rangle.
\]
Dla pewnego $m\in\N_0$ i każdego wielowskaźnika $\beta$
istnieje $C_\beta>0$ takie, że
$|\partial_\xi^\beta F_T(\xi)|
\le C_\beta(1+|\xi|)^m$
dla wszystkich $\xi\in\R^d$.
W działaniu na funkcji wykładniczej używamy
przedłużenia z Lematu~\ref{lem:compact-action21}.
\end{twierdzenie}
\begin{proof}
Wybierzmy $\chi\in\mathcal D$ równą $1$ na
otoczeniu $K$. Dla $\beta\in\N_0^d$ funkcje
$x\mapsto\chi(x)(-ix)^\beta e^{-ix\cdot\xi}$
mają nośniki w ustalonym zwartym zbiorze $\supp\chi$.
Na zwartym otoczeniu ustalonego parametru $\xi$
wszystkie pochodne względem $x$ i parametru są ciągłe
i jednostajnie ciągłe na iloczynie tych dwóch zbiorów.
Stosując twierdzenie Newtona--Leibniza do ilorazów
różnicowych względem parametru w seminormach
$\mathcal D_{\supp\chi}$, dostajemy
\[
 \partial_\xi^\beta F_T(\xi)
 =(2\pi)^{-d/2}
 \langle T_x,(-ix)^\beta e^{-ix\cdot\xi}\rangle.
\]
Ciągłość tych pochodnych wynika z tego samego argumentu.
Z Lematu~\ref{lem:compact-action21} wybieramy $C,m,L$
niezależne od $\beta,\xi$, gdzie $L$ jest zwarty i
$|T(f)|\le C\max_{|\alpha|\le m}\sup_{x\in L}
|\partial^\alpha f(x)|$.
Dla $|\alpha|\le m$ wzór Leibniza daje
\[
 \partial_x^\alpha\bigl((-ix)^\beta e^{-ix\cdot\xi}\bigr)
 =\sum_{\gamma\le\alpha}\binom\alpha\gamma
 \partial_x^\gamma((-ix)^\beta)
 (-i\xi)^{\alpha-\gamma}e^{-ix\cdot\xi}.
\]
Ponieważ $x\in L$, pochodne wielomianu są ograniczone,
a moduł czynnika wykładniczego dla rzeczywistych $\xi$
wynosi $1$. Ponadto $|\xi^{\alpha-\gamma}|
\le(1+|\xi|)^m$. Po zsumowaniu skończonej liczby
składników otrzymujemy $C_\beta(1+|\xi|)^m$ jako
oszacowanie modułu $\partial_\xi^\beta F_T(\xi)$.
W szczególności $T_{F_T}$ jest temperowana.

Aby sprawdzić równość dystrybucji, ustalmy
$\varphi\in\mathcal S$ i $c_d=(2\pi)^{-d/2}$.
Dla $R\ge1$ połóżmy
\[
 J_R(x)=c_d\chi(x)\int_{[-R,R]^d}
 e^{-ix\cdot\xi}\varphi(\xi)\,\mathrm d\lambda_d(\xi).
\]
Sumy Riemanna zbiegają do $J_R$ w
$\mathcal D_{\supp\chi}$: dla każdego rzędu pochodnych
względem $x$ wynika to z jednostajnej ciągłości funkcji
podcałkowej i jej pochodnych na iloczynie zwartych
zbiorów $\supp\chi$ i $[-R,R]^d$.
Stąd
\[
 \langle T,J_R\rangle=
 \int_{[-R,R]^d}F_T(\xi)\varphi(\xi)
 \,\mathrm d\lambda_d(\xi).
\]
Dla każdego $q\in\N_0$ wzór Leibniza daje
\[
 p_q(J_R-\chi\mathcal F\varphi)
 \le C_{\chi,q,d}\int_{\R^d\setminus[-R,R]^d}
 (1+|\xi|)^q|\varphi(\xi)|\,\mathrm d\lambda_d(\xi)
 \longrightarrow0.
\]
Pochodne pod całką istnieją na mocy całkowalności
każdej wagi wielomianowej razy $|\varphi|$.
Wszystkie funkcje po lewej mają nośniki w $\supp\chi$,
więc zbieżność zachodzi w jednym etapie $\mathcal D$.
Z ciągłości $T$ i bezwzględnej całkowalności
$F_T\varphi$ dostajemy
\[
 \langle\mathcal FT,\varphi\rangle
 =\langle T,\chi\mathcal F\varphi\rangle
 =\int_{\R^d}F_T(\xi)\varphi(\xi)
 \,\mathrm d\lambda_d(\xi).
\]
Pierwsza równość wykorzystuje jednoznaczne przedłużenie
$T$ na $\mathcal S$ i regułę odcięcia.
Zatem $\mathcal FT=T_{F_T}$.
\end{proof}

\begin{uwaga}[Zakres rozdziału]
\label{rem:scope22}
Transformacja Fouriera jest tu automorfizmem przestrzeni
dystrybucji temperowanych i podstawowym narzędziem teorii
operatorów. Jej zastosowania do rozwiązywania konkretnych
równań różniczkowych pozostawiamy osobnemu skryptowi.
Dalsze rozdziały rozwiną teorię przestrzeni Montela,
przestrzeni nuklearnych i produktów tensorowych topologicznych.
\end{uwaga}

% Materiał z AF 25; fragment do rozdziału o dystrybucjach temperowanych.
\subsection{Łańcuch przestrzeni sprzężonych i kryterium temperowania}

\begin{lemat}[Transpozycja gęstego zanurzenia]
\label{lem:transpozycja25}
Niech $X,Y$ będą przestrzeniami lokalnie wypukłymi
Hausdorffa, a $j:X\to Y$ ciągłym odwzorowaniem liniowym
o gęstym obrazie. Wówczas transpozycja
$j':Y'\to X'$, $j'(T)=T\circ j$, jest liniowa,
różnowartościowa i ciągła zarówno dla topologii słabych-*, jak i dla
silnych topologii przestrzeni dualnych.
\end{lemat}
\begin{proof}
Dla każdego $T\in Y'$ złożenie $T\circ j$ jest ciągłym
funkcjonałem liniowym na $X$, więc $j'$ jest dobrze
określone; liniowość wynika z liniowości składania
względem pierwszego czynnika. Jeżeli $T\circ j=0$,
to $T=0$ przez gęstość $j(X)$ i ciągłość $T$.
Dla $x\in X$ mamy
$|\langle j'T,x\rangle|=|\langle T,jx\rangle|$,
co daje ciągłość w topologiach słabych-*.
Jeśli $B\subseteq X$ jest ograniczony, to $j(B)$ jest
ograniczony w $Y$ (dla każdej ciągłej seminormy $p$
na $Y$ seminorma $p\circ j$ jest ciągła na $X$).
Dla seminorm silnych zachodzi dokładna równość
$\sup_{x\in B}|\langle j'T,x\rangle|
= \sup_{y\in j(B)}|\langle T,y\rangle|$.
To dowodzi ciągłości w topologiach silnych.
\end{proof}

\begin{wniosek}[Łańcuch przestrzeni dualnych]
\label{cor:dual-lancuch25}
Dla każdego $d\ge1$ zachodzą kanoniczne, ciągłe i
różnowartościowe odwzorowania liniowe
\[
 \mathcal E'(\R^d)\longrightarrow\mathcal S'(\R^d)
 \longrightarrow\mathcal D'(\R^d),
\]
zarówno w topologiach słabych-*, jak i silnych.
Dla każdego niepustego zbioru otwartego $\Omega\subseteq\R^d$
kanoniczne odwzorowanie
$\mathcal E'(\Omega)\to\mathcal D'(\Omega)$
ma te same własności. Samo stwierdzenie ciągłości i
różnowartościowości \emph{nie oznacza}, że są to zanurzenia
topologiczne na obrazy z topologią podprzestrzeni.
\end{wniosek}
\begin{proof}
Stosujemy Lemat~\ref{lem:transpozycja25} do odwzorowań
z Twierdzenia~\ref{thm:zanurzenia25}. Jeżeli $i:X\to Y$
oraz $j:Y\to Z$ są tymi kolejnymi odwzorowaniami,
to $(j\circ i)'=i'\circ j'$, gdyż dla $T\in Z'$
mamy $(T\circ j)\circ i=T\circ(j\circ i)$.
Dlatego ograniczenie z $\mathcal E'$ do $\mathcal D'$
przez $\mathcal S'$ jest tym samym odwzorowaniem co
bezpośrednie ograniczenie. W przypadku zespolonym
transpozycja jest liniowa, bez sprzężenia zespolonego.
Ciągłość i różnowartościowość nie zapewniają ciągłości
odwrotności na obrazie; poniżej podamy kontrprzykłady.
\end{proof}
\subsection{Dystrybucje temperowane jako podprzestrzeń \texorpdfstring{$\mathcal D'$}{D'}}

\begin{twierdzenie}[Kryterium temperowania]
\label{thm:temp-criterion25}
Niech $d\ge1$ będzie liczbą całkowitą, $\K\in\{\R,\C\}$
i $T\in\mathcal D'(\R^d)$ nad $\K$.
Następujące warunki są równoważne:
\begin{enumerate}[label=\textup{(\roman*)}]
\item $T$ należy do obrazu kanonicznego odwzorowania
$\mathcal S'(\R^d)\to\mathcal D'(\R^d)$;
\item istnieją $C>0$, $N,m\in\N_0$ takie, że
$|\langle T,\varphi\rangle|\le C p_{N,m}(\varphi)$
dla każdej $\varphi\in\mathcal D(\R^d)$.
\end{enumerate}
Jeżeli te warunki zachodzą, przedłużenie
$\widetilde T\in\mathcal S'(\R^d)$ jest jednoznaczne
i spełnia $|\langle\widetilde T,f\rangle|\le C p_{N,m}(f)$
dla wszystkich $f\in\mathcal S(\R^d)$, z tymi samymi
stałymi co w~\textup{(ii)}. W szczególności każda
dystrybucja o zwartym nośniku jest temperowana.
\end{twierdzenie}
\begin{proof}
Niech najpierw $\widetilde T\in\mathcal S'$ będzie
przedłużeniem $T$. Z ciągłości w zerze i współkońcowości
rosnącego ciągu norm $p_{j,j}$ istnieją $j\in\N_0$ oraz
$\delta>0$ takie, że $p_{j,j}(f)<\delta$ implikuje
$|\widetilde T(f)|<1$. Dla $f\ne0$ stosujemy ten warunek
do $\delta f/(2p_{j,j}(f))$, otrzymując
$|\widetilde T(f)|\le (2/\delta)p_{j,j}(f)$.
Dla $f=0$ nierówność jest oczywista. Ograniczenie
do $\mathcal D$ daje warunek~\textup{(ii)}.
Jest to argument z Twierdzenia~\ref{thm:tempered-estimate22},
działający nad obydwoma ciałami.

Załóżmy teraz~\textup{(ii)}. Dla każdego $f\in\mathcal S$
wybieramy ciąg $(\varphi_j)$ w $\mathcal D$ zbieżny do $f$
w $\mathcal S$ (Twierdzenie~\ref{thm:zanurzenia25}).
Ponieważ
\[
 |T(\varphi_j)-T(\varphi_k)|
 \le C p_{N,m}(\varphi_j-\varphi_k)\longrightarrow0
 \quad(j,k\to\infty),
\]
ciąg wartości ma granicę w zupełnym ciele $\K$.
Kładziemy $\widetilde T(f)=\lim_jT(\varphi_j)$.
Dla drugiego ciągu $(\psi_j)$ zbieżnego do $f$ mamy
$|T(\varphi_j)-T(\psi_j)|\le
C p_{N,m}(\varphi_j-\psi_j)\to0$, więc definicja nie zależy
od wyboru ciągu. Dla $f,g\in\mathcal S$ i $a,b\in\K$
można aproksymować $af+bg$ przez $a\varphi_j+b\psi_j$,
gdzie teraz $\varphi_j\to f$ i $\psi_j\to g$.
Przechodząc do granicy, otrzymujemy liniowość.
Ponadto $|p_{N,m}(\varphi_j)-p_{N,m}(f)|
\le p_{N,m}(\varphi_j-f)$, więc
\[
 |\widetilde T(f)|\le C p_{N,m}(f).
\]
Daje to ciągłość; na $\mathcal D$ wystarczy wziąć
ciąg stały, aby uzyskać $\widetilde T|_{\mathcal D}=T$.
Jednoznaczność wynika z gęstości $\mathcal D$.

Jeżeli $T$ ma zwarty nośnik, to z
Twierdzenia~\ref{thm:compact-support-E25} otrzymujemy
$|T(\varphi)|\le C p_{K,m}(\varphi)\le C p_{0,m}(\varphi)$
na funkcjach testowych, a zatem~\textup{(ii)} zachodzi z $N=0$.
Przedłużenie jest ograniczeniem funkcjonału na $\mathcal E$
do $\mathcal S$, przez jednoznaczność; odpowiada to
Stwierdzeniu~\ref{prop:schwartz-embeddings22}.
\end{proof}

\begin{przyklad}[Dystrybucja nietemperowana]
\label{ex:non-tempered25}
Na $\R$ funkcja $f(x)=e^x$ jest lokalnie całkowalna,
więc wyznacza dystrybucję regularną $T_f\in\mathcal D'(\R)$.
Nie jest ona temperowana. Wybierzmy nieujemną, niezerową
$\psi\in\mathcal D((-1,1))$, przedłużoną zerem na $\R$,
i niech $\varphi_n(x)=e^{-n/2}\psi(x-n)$ dla $n\ge1$.
Dla $x\in\supp\varphi_n$ mamy $|x|\le n+1$; pochodne
$\varphi_n$ znikają poza tym nośnikiem. Stąd dla wszystkich
$N,m\in\N_0$
\[
 p_{N,m}(\varphi_n)\le e^{-n/2}(n+2)^N p_{0,m}(\psi)
 \longrightarrow0,
\]
gdyż funkcja wykładnicza rośnie szybciej niż każdy
wielomian. Jednak po podstawieniu $y=x-n$ otrzymujemy
\[
 \langle T_f,\varphi_n\rangle=
 e^{n/2}\int_\R e^y\psi(y)\,\mathrm d\lambda_1(y)
 \longrightarrow+\infty.
\]
Całka jest skończona i ściśle dodatnia: $\psi$ jest
nieujemna i dodatnia na pewnym niepustym przedziale.
Ciągłość przedłużenia na $\mathcal S$ wymagałaby
$\langle T_f,\varphi_n\rangle\to0$, co daje sprzeczność.
Przykład działa nad $\R$ i $\C$; razem z
Przykładem~\ref{ex:non-tempered22} pokazuje on dwie
metody wykrywania braku temperowalności (przesuwanie
nośnika oraz powiększanie odcięcia).
\end{przyklad}
\begin{przyklad}[Zbieżność i ucieczka nośnika]
\label{ex:escaping-deltas25}
Na $\R$ niech $\delta_n$ oznacza deltę Diraca w punkcie
$n$ dla $n\ge1$. Każda z tych dystrybucji ma zwarty
nośnik, więc należy także do $\mathcal E'$ i $\mathcal S'$
na mocy Twierdzeń~\ref{thm:compact-support-E25}
i~\ref{thm:temp-criterion25}.
Wtedy $\delta_n\to0$ silnie w $\mathcal D'(\R)$:
dla ograniczonego $B\subseteq\mathcal D(\R)$ nośniki
wszystkich jego elementów mieszczą się w jednym zwartym
przedziale (Twierdzenie~\ref{thm:D-bounded18}), więc
$\sup_{\varphi\in B}|\varphi(n)|=0$ dla dużych $n$.
Również $\delta_n\to0$ silnie w $\mathcal S'(\R)$, bo dla
ograniczonego zbioru $B\subseteq\mathcal S$ mamy
$C_B=\sup_{\varphi\in B}p_{2,0}(\varphi)<\infty$ oraz
$\sup_{\varphi\in B}|\varphi(n)|\le C_B(1+n)^{-2}$.
Jednak $\delta_n\not\to0$ nawet słabo-* w
$\mathcal E'(\R)$, ponieważ $\langle\delta_n,1\rangle=1$.

Z kolei $T_n=e^n\delta_n$ zbiega silnie do zera w
$\mathcal D'(\R)$ (na każdym ograniczonym zbiorze
seminorma silna jest od pewnego indeksu zerowa).
Nie zbiega do zera nawet słabo-* w $\mathcal S'(\R)$.
Niech $r(x)=\sqrt{1+x^2}$ i $\varphi(x)=e^{-r(x)/2}$.
Wykażemy, że $\varphi\in\mathcal S(\R)$.
Dla każdego $k\in\N_0$ istnieje wielomian $P_k(x,z)$
taki, że
\[
 \varphi^{(k)}(x)=e^{-r(x)/2}r(x)^{-2k}P_k(x,r(x)).
\]
Dla $k=0$ przyjmujemy $P_0=1$. Ponieważ $r'=x/r$,
różniczkowanie wzoru daje rekurencję
\[
 P_{k+1}=z^2\partial_x P_k+xz\partial_z P_k
          -2kxP_k-\tfrac12 xzP_k.
\]
Zatem wszystkie $P_k$ są wielomianami. Ponieważ
$1\le r(x)\le1+|x|$ i $r(x)\ge|x|$, każda pochodna
jest ograniczona przez stałą razy wielomianowy wzrost
$(1+|x|)^{M_k}$ pomnożony przez $e^{-|x|/2}$.
Wszystkie seminormy Schwartza są więc skończone.
Jednak dla $n\ge1$ mamy $r(n)\le n+1$, toteż
$\langle T_n,\varphi\rangle=e^{n-r(n)/2}
\ge e^{(n-1)/2}\to+\infty$.

Dla dowolnego niepustego zbioru otwartego $\Omega$
wybierzmy wyczerpanie $(K_j)$ i punkty
$a_j\in\Omega\setminus K_j$.
Wówczas $\delta_{a_j}\to0$ silnie w $\mathcal D'(\Omega)$,
gdyż każdy ograniczony zbiór testów ma wspólny zwarty
nośnik zawarty w pewnym $K_J$.
Natomiast $\langle\delta_{a_j},1\rangle=1$ na
$\mathcal E(\Omega)$ dla wszystkich $j$.
Zatem wszystkie odwzorowania z
Wniosku~\ref{cor:dual-lancuch25} nie są izomorfizmami
topologicznymi na swoje obrazy z topologią podprzestrzeni,
ani dla topologii słabych-*, ani dla silnych.
\end{przyklad}

\subsection{Zgodność operacji i granice ich działania}

\begin{stwierdzenie}[Różniczkowanie i mnożenie]
\label{prop:ops25}
Niech $\Omega\subseteq\R^d$ będzie niepustym zbiorem otwartym.
Dla każdego wielowskaźnika $\alpha\in\N_0^d$
operator $\partial^\alpha$ działa liniowo i ciągle na
$\mathcal E(\Omega)$ i $\mathcal D(\Omega)$, a przez
transpozycję także na $\mathcal E'(\Omega)$ i
$\mathcal D'(\Omega)$ w topologiach słabych-* i silnych.
Dla $a\in C^\infty(\Omega)$ mnożenie $f\mapsto af$
i odpowiadające mu mnożenie dystrybucji działają ciągle
na tych przestrzeniach.
Różniczkowanie działa liniowo i ciągle również na
$\mathcal S(\R^d)$ i na $\mathcal S'(\R^d)$ w obu
topologiach dualnych. Na tych przestrzeniach
te same własności mnożenia zachodzą dla
$a\in C^\infty(\R^d)$, jeżeli dla każdego
wielowskaźnika $\alpha$ istnieją $C_\alpha>0$ i
$N_\alpha\in\N_0$ takie, że
$|\partial^\alpha a(x)|\le C_\alpha(1+|x|)^{N_\alpha}$
dla wszystkich $x\in\R^d$.
\end{stwierdzenie}
\begin{proof}
Dla $\mathcal E$ i $\mathcal D$ oszacowania seminorm
dla różniczkowania i mnożenia wynikają z reguły Leibniza,
jak w stwierdzeniu~\ref{prop:operations-cinf17}
i twierdzeniach~\ref{thm:dist-deriv-cont20}, \ref{thm:multiply20}.
Dla $\mathcal S$ i $\mathcal S'$ stosujemy stwierdzenie
\ref{prop:schwartz-operators17} oraz twierdzenie
\ref{thm:tempered-operators22}, w tym oszacowanie
mnożenia przez funkcję o pochodnych wzrostu wielomianowego.
Dla ciągłego operatora $A:X\to X$ transpozycja spełnia
$P_B(A'T)=P_{A(B)}(T)$, gdzie
$P_B(T)=\sup_{f\in B}|T(f)|$ i $A(B)$ jest ograniczony.
Daje to ciągłość silną, a dla singletonów także słabą-*.
Stosujemy ten argument w szczególności do
$X=\mathcal E$ i operatorów $(-1)^{|\alpha|}\partial^\alpha$
oraz $f\mapsto af$. Operatory na funkcjach komutują
z kanonicznymi odwzorowaniami przestrzeni funkcji,
więc przez transpozycję odpowiadające operacje
na przestrzeniach sprzężonych komutują z ograniczeniami
z wniosku~\ref{cor:dual-lancuch25} (w przypadku $\mathcal S'$
przy podanym warunku wzrostu mnożnika).
\end{proof}

\begin{przyklad}[Nie każdy mnożnik gładki zachowuje przestrzeń Schwartza]
\label{ex:not-multiplier25}
Funkcja $a(x)=e^{|x|^2}$ jest gładka, ale mnożenie przez nią
nie zachowuje przestrzeni Schwartza: dla
$f(x)=e^{-|x|^2}\in\mathcal S(\R^d)$ iloczyn $af=1$
nie należy do $\mathcal S(\R^d)$. Przynależność funkcji Gaussa do $\mathcal S$ wynika
z tego, że każda jej pochodna jest wielomianem
pomnożonym przez tę funkcję (indukcja z reguły
iloczynu), a spadek Gaussa dominuje każdą potęgę.
Funkcja $1$ nie spełnia warunku $p_{1,0}(1)<\infty$.
W przestrzeniach $\mathcal D$ i $\mathcal E$
mnożenie przez $a$ jest natomiast ciągłe.
Co więcej, zwykła dystrybucja $aT_1=T_a$ nie jest
temperowana, choć $T_1$ jest temperowana
(Przykłady~\ref{ex:tempered-regular22}
i~\ref{ex:non-tempered22}).
\end{przyklad}

\begin{uwaga}[Transformacja Fouriera i splot]
\label{rem:fourier-conv25}
Nad $\C$ transformacja Fouriera jest automorfizmem
liniowo-topologicznym $\mathcal S(\R^d)$ i
$\mathcal S'(\R^d)$ w obu topologiach dualnych
(Twierdzenia~\ref{thm:fourier-schwartz17}
i~\ref{thm:fourier-auto-tempered22}).
Nie jest zdefiniowana przez tę samą transpozycję
na całej $\mathcal D'$: niezerowa funkcja testowa
nie może mieć transformaty Fouriera o zwartym nośniku.
Istotnie, dla $f\in\mathcal D$ jej transformata
jest rzeczywiście analityczna: rozwinięcie wykładnika
w szereg potęgowy w otoczeniu dowolnego punktu
można całkować wyraz po wyrazie, ponieważ $f$ ma
zwarty nośnik, a szereg jest jednostajnie bezwzględnie
zbieżny na nim dla parametrów z każdego ograniczonego
otoczenia. Gdyby transformata miała zwarty nośnik,
znikałaby na niepustym otwartym podzbiorze $\R^d$,
a z jednoznaczności funkcji rzeczywiście analitycznej
znikałaby na całym spójnym $\R^d$.
Różnowartościowość transformaty na $\mathcal S$
dawałaby $f=0$. Zatem wyrażenie
$\langle T,\mathcal F f\rangle$ nie musi mieć sensu
dla $T\in\mathcal D'$ i $f\in\mathcal D$.

Dla ustalonego $T\in\mathcal E'(\R^d)$ splot
$S\mapsto T*S$ jest ciągłym operatorem liniowym
na $\mathcal D'(\R^d)$ w obu topologiach
(Twierdzenie~\ref{thm:convolution-properties21}).
Aby sprawdzić, że zachowuje również $\mathcal S'$,
niech $L$ będzie zwartym zbiorem i $C>0$, $m\in\N_0$
stałymi oszacowania przedłużenia $T$ na $\mathcal E$
z Twierdzenia~\ref{thm:compact-support-E25}.
Dla $f\in\mathcal S$ określamy
$(A_Tf)(y)=\langle T_x,f(x+y)\rangle$,
korzystając z przedłużenia na $\mathcal E$.
Z ciągłości $T$ na $\mathcal E$ i różniczkowania
przesunięć w jej topologii otrzymujemy
$\partial_y^\beta A_T f(y)=
\langle T_x,\partial^\beta f(x+y)\rangle$:
ilorazy różnicowe wraz z pochodnymi względem $x$
zbiegają jednostajnie na każdym zwartym zbiorze
przez wzór całkowy na różnicę funkcji gładkiej.
Ponieważ $1+|y|\le(1+|x|)(1+|x+y|)$, dla
$R=\max_{x\in L}|x|$ otrzymujemy
\[
 p_{N,k}(A_Tf)\le C(1+R)^N p_{N,k+m}(f).
\]
Zatem $A_T:\mathcal S\to\mathcal S$ jest liniowe
i ciągłe. Jego transpozycja
$\langle T*S,f\rangle=\langle S,A_Tf\rangle$
jest ciągła na $\mathcal S'$ w obu topologiach
i zgadza się na $\mathcal D$ z definicją splotu
z rozdziału~\ref{sec:splot21}.
Nie wynika stąd istnienie splotu dwóch dowolnych
dystrybucji (Uwaga~\ref{rem:convolution-limits21}).
\end{uwaga}

\subsection{Transformacja całkowa w przestrzeni \texorpdfstring{$L^1$}{L1}}

Przyjmujemy w całym skrypcie normalizację unitarną z
Definicji~\ref{def:fourier-schwartz17}. Rozróżniamy
transformację (operator) i transformatę (jego wartość).
W $L^1$ całka określa funkcję w każdym punkcie, natomiast
w $L^2$ twierdzenie Plancherela określa w pierwszej
kolejności klasę funkcji równych prawie wszędzie.
Dla $f\in L^1(\R^d)$ definiujemy
\[
 \widehat f(\xi)=c_d\int_{\R^d}f(x)e^{-ix\cdot\xi}
 \,\mathrm d\lambda_d(x),\qquad c_d=(2\pi)^{-d/2}.
\]
Intuicyjnie jest to pomiar amplitudy i fazy poszczególnych
częstotliwości; przy dużej częstotliwości wkłady z różnych
punktów mogą się znosić, nawet gdy moduł jądra jest równy
jeden. To zjawisko opisuje następne twierdzenie.

\begin{twierdzenie}[Lemat Riemanna--Lebesgue'a]
\label{thm:dt-fourier-l1-rl}
Transformacja całkowa jest ograniczonym operatorem
liniowym z $L^1(\R^d)$ do $C_0(\R^d)$, gdzie $C_0$
oznacza przestrzeń ciągłych funkcji zanikających w
nieskończoności, z normą supremum. Zachodzi
$\|\widehat f\|_\infty\le c_d\|f\|_1$.
\end{twierdzenie}
\begin{proof}
Liniowość i oszacowanie wynikają bezpośrednio z całki.
Ciągłość jest konsekwencją twierdzenia o zbieżności
zmajoryzowanej z majorantą $|f|$.
Aby wykazać zanikanie, dla $h\in\R^d$ oznaczmy
$\tau_h f(x)=f(x-h)$. Z niezmienniczości miary
Lebesgue'a względem translacji
$\widehat{\tau_h f}(\xi)=e^{-ih\cdot\xi}\widehat f(\xi)$.
Dla $\xi\ne0$ wybieramy $h=\pi\xi/|\xi|^2$.
Wtedy $e^{-ih\cdot\xi}=-1$ oraz
\[
 2|\widehat f(\xi)|
 =|\widehat{f-\tau_h f}(\xi)|
 \le c_d\|f-\tau_h f\|_1.
\]
Gdy $|\xi|\to\infty$, mamy $h\to0$, a prawa strona
dąży do zera na mocy
\atlasref{ME-14::thm:translation-lp-14}{ciągłości translacji w przestrzeni $L^1$}.
\end{proof}

\begin{stwierdzenie}[Reguły transformacji całkowej]
\label{prop:dt-fourier-l1-rules}
Dla $f,g\in L^1(\R^d)$, $h,\eta\in\R^d$ i odwracalnej
rzeczywistej macierzy $A$ zachodzą tożsamości
\[
\begin{aligned}
 \widehat{\tau_h f}(\xi)&=e^{-ih\cdot\xi}\widehat f(\xi),\\
 \mathcal F(e^{i\eta\cdot x}f(x))(\xi)&=\widehat f(\xi-\eta),\\
 \mathcal F(f(A\,\cdot))(\xi)&=
 |\det A|^{-1}\widehat f(A^{-T}\xi),\\
 \widehat{f*g}(\xi)&=c_d^{-1}\widehat f(\xi)\widehat g(\xi).
\end{aligned}
\]
Jeżeli ponadto $x^\beta f\in L^1$ dla wszystkich
wielowskaźników $|\beta|\le m$, to $\widehat f\in C^m$
oraz $\partial_\xi^\beta\widehat f=(-i)^{|\beta|}
\mathcal F(x^\beta f)$.
\end{stwierdzenie}
\begin{proof}
Pierwsza tożsamość wystąpiła w dowodzie
Twierdzenia~\ref{thm:dt-fourier-l1-rl}; druga wynika z
połączenia wykładników. W trzeciej stosujemy podstawienie
$y=Ax$, dla którego $x\cdot\xi=y\cdot A^{-T}\xi$.
W ostatniej całkowalność modułu funkcji podcałkowej
wynika z
\[
 \int_{\R^d}\int_{\R^d}|f(y)g(x-y)|
 \,\mathrm d\lambda_d(y)\,\mathrm d\lambda_d(x)
 =\|f\|_1\|g\|_1.
\]
Twierdzenie Fubiniego i podstawienie $z=x-y$ pozwalają
rozłożyć całkę na iloczyn dwóch całek. Czynnik wynosi
$c_d/c_d^2=c_d^{-1}$.
Dla pochodnych różniczkujemy pod całką, korzystając z
$|(e^{-itx_j}-1)/t|\le |x_j|$. Na każdym kolejnym etapie
majorantą jest $|x^{\beta+e_j}f(x)|$; ciągłość
pochodnych wynika ponownie ze zbieżności zmajoryzowanej.
\end{proof}

\begin{stwierdzenie}[Zgodność transformat całkowej, dystrybucyjnej i unitarnej]
\label{prop:dt-fourier-l1-l2-compat}
Dla $f\in L^1$ transformacja dystrybucyjna spełnia
$\mathcal FT_f=T_{\widehat f}$, gdzie $\widehat f$
jest transformatą całkową. Gdy $f\in L^1\cap L^2$,
transformata całkowa jest reprezentantem transformaty
unitarnej z Twierdzenia~\ref{thm:plancherel22}.
\end{stwierdzenie}
\begin{proof}
Dla $\varphi\in\mathcal S$ całkowalność modułu
$|f(x)\varphi(\xi)|$ wynika z $f,\varphi\in L^1$.
Z twierdzenia Fubiniego otrzymujemy
\[
 \int_{\R^d}\widehat f(\xi)\varphi(\xi)\,\mathrm d\lambda_d(\xi)
 =\int_{\R^d}f(x)(\mathcal F\varphi)(x)\,\mathrm d\lambda_d(x).
\]
Jest to definicyjna tożsamość transformacji dystrybucyjnej.
Dla $f\in L^1\cap L^2$ twierdzenie Plancherela daje
tę samą dystrybucję regularną dla transformaty w $L^2$.
Oba reprezentanty są lokalnie całkowalne; ich równość
prawie wszędzie wynika ze
Stwierdzenia~\ref{prop:regular-injective19}.
\end{proof}

\begin{lemat}[Odwracanie z regularyzacją Gaussa]
\label{lem:dt-fourier-gaussian-l1}
Niech $f\in L^1(\R^d)$ i $\varepsilon>0$.
Definiujemy jądro Gaussa
\[
 G_\varepsilon(x)=(2\pi\varepsilon)^{-d/2}
 e^{-|x|^2/(2\varepsilon)}.
\]
Wtedy dla każdego $x$ zachodzi
\[
 c_d\int_{\R^d}e^{ix\cdot\xi}
 e^{-\varepsilon|\xi|^2/2}\widehat f(\xi)
 \,\mathrm d\lambda_d(\xi)=(G_\varepsilon*f)(x).
\]
Ponadto $G_\varepsilon*f\to f$ w $L^1$ dla
$\varepsilon\downarrow0$.
\end{lemat}
\begin{proof}
Lewa strona jest bezwzględnie zbieżna na mocy
Twierdzenia~\ref{thm:dt-fourier-l1-rl}.
Podstawiamy definicję $\widehat f$ i stosujemy Fubiniego,
gdyż $|f(y)|e^{-\varepsilon|\xi|^2/2}$ jest całkowalne
na przestrzeni produktowej. Całkę po $\xi$ obliczamy
wzorem Gaussa z dowodu Twierdzenia~\ref{thm:fourier-schwartz17};
łącznie z czynnikiem $c_d^2$ daje to $G_\varepsilon(x-y)$.
Z tego samego dowodu wiemy, że $\int G_\varepsilon=1$
oraz całka $G_\varepsilon$ poza dowolną kulą o dodatnim
promieniu dąży do zera. Twierdzenie Tonellego daje
\[
 \|G_\varepsilon*f-f\|_1
 \le \int_{\R^d}G_\varepsilon(y)
 \|\tau_yf-f\|_1\,\mathrm d\lambda_d(y).
\]
Dla dowolnego $\delta>0$ część po $|y|<\delta$ nie
przekracza $\sup_{|y|<\delta}\|\tau_yf-f\|_1$,
a pozostała nie przekracza
$2\|f\|_1\int_{|y|\ge\delta}G_\varepsilon\,\mathrm d\lambda_d$.
Ciągłość translacji pozwala najpierw wybrać $\delta$,
a następnie zmniejszyć $\varepsilon$, tak aby obie
części były dowolnie małe.
\end{proof}

\begin{twierdzenie}[Odwrócenie i jednoznaczność w $L^1$]
\label{thm:dt-fourier-l1-inversion}
Jeżeli $f,\widehat f\in L^1(\R^d)$, to
\[
 f(x)=c_d\int_{\R^d}e^{ix\cdot\xi}\widehat f(\xi)
 \,\mathrm d\lambda_d(\xi)
\]
dla prawie każdego $x$, a prawa strona jest ciągłym,
zanikającym w nieskończoności reprezentantem $f$.
Jeżeli $f$ jest ciągła, równość zachodzi wszędzie.
Niezależnie od całkowalności $\widehat f$ transformacja
na $L^1$ jest różnowartościowa.
\end{twierdzenie}
\begin{proof}
Niech $g$ oznacza prawą stronę wzoru odwrotnego.
Twierdzenie~\ref{thm:dt-fourier-l1-rl}, ze zmienionym
znakiem argumentu, daje $g\in C_0$.
Z Lematu~\ref{lem:dt-fourier-gaussian-l1} oraz
zbieżności zmajoryzowanej mamy
\[
 \|G_\varepsilon*f-g\|_\infty
 \le c_d\int_{\R^d}
 (1-e^{-\varepsilon|\xi|^2/2})|\widehat f(\xi)|
 \,\mathrm d\lambda_d(\xi)\longrightarrow0.
\]
Z drugiej strony $G_\varepsilon*f\to f$ w $L^1$.
Na każdej kuli skończonej miary zbieżność jednostajna
implikuje zbieżność w $L^1$ na tej kuli. Jednoznaczność
granicy daje $g=f$ prawie wszędzie na każdej takiej kuli,
a ich przeliczalne wyczerpanie daje równość na $\R^d$.
Jeżeli $f$ jest ciągła, różnica dwóch funkcji ciągłych
równa zeru prawie wszędzie jest identycznie zerowa.
Jeżeli $\widehat f=0$, nawet bez założenia o
całkowalności transformaty, lewa strona tożsamości
z Lematu~\ref{lem:dt-fourier-gaussian-l1} jest zerowa
dla każdego $\varepsilon$, więc $G_\varepsilon*f=0$.
Przejście do granicy w $L^1$ daje $f=0$.
Liniowość implikuje różnowartościowość.
\end{proof}

\begin{przyklad}[Prostokąt i funkcja wykładnicza]
\label{ex:dt-fourier-l1-pairs}
Dla $a>0$ w wymiarze jeden mamy
\[
 \mathcal F\chi_{[-a,a]}(\xi)
 =\sqrt{\frac{2}{\pi}}\frac{\sin(a\xi)}{\xi},
 \qquad
 \mathcal F(e^{-a|x|})(\xi)
 =\sqrt{\frac{2}{\pi}}\frac{a}{a^2+\xi^2}.
\]
Pierwszą wartość przy $\xi=0$ rozumiemy jako
$\sqrt{2/\pi}\,a$. Pierwszy wzór wynika z całkowania
$e^{-ix\xi}$ po $[-a,a]$. W drugim rozdzielamy całkę
na półproste, otrzymując
$c_1((a+i\xi)^{-1}+(a-i\xi)^{-1})$.
Transformata prostokąta nie należy do $L^1$: na
przedziałach $a\xi\in[n\pi+\pi/6,n\pi+5\pi/6]$
dla $n\ge1$ licznik ma moduł co najmniej $1/2$,
więc suma całek jest ograniczona z dołu przez
dodatnią stałą razy $\sum_{n\ge1}(n+1)^{-1}$.
W drugim przypadku obie funkcje należą do $L^1$,
a wzór odwrotny obowiązuje w każdym punkcie.
\end{przyklad}

\begin{uwaga}[Zgodność z normalizacją RRC]
\label{rem:dt-fourier-normalization-rrc}
W skrypcie RRC-08 stosowana jest transformata
$\mathcal F_0f(\xi)=\int e^{-ix\cdot\xi}f(x)
\,\mathrm d\lambda_d(x)$. Mamy $\mathcal F=c_d\mathcal F_0$;
we wzorze odwrotnym dla $\mathcal F_0$ występuje
$(2\pi)^{-d}$, a we wzorze na splot nie ma dodatkowego
czynnika. To nie są odmienne twierdzenia, lecz różne
normalizacje tego samego operatora. Por.
\atlasref{PDE-08::prop:pde8-RL}{RRC-08: lemat Riemanna--Lebesgue'a}
oraz
\atlasref{PDE-08::thm:pde8-inversion}{RRC-08: odwrócenie transformaty}.
\end{uwaga}

\subsection{Interpolacja zespolona i nierówność Hausdorffa--Younga}

Wynik dla $L^1$ daje kontrolę normy supremum transformaty,
a twierdzenie Plancherela zachowanie normy $L^2$.
Pomiędzy tymi przypadkami można uzyskać oszacowania w
$L^{p'}$, gdzie $p'$ jest wykładnikiem sprzężonym do $p$.
Potrzebujemy jednak interpolacji norm operatorów, nie tylko
interpolacji norm jednej funkcji.

\begin{lemat}[Lemat o trzech prostych]
\label{lem:dt-three-lines}
Niech $F$ będzie ograniczoną funkcją ciągłą na domkniętym
pasie $0\le\operatorname{Re}z\le1$, holomorficzną wewnątrz
pasa. Załóżmy, że dla każdego $t\in\R$ zachodzą
$|F(it)|\le M_0$ oraz $|F(1+it)|\le M_1$.
Dla $0<\theta<1$ mamy
\[
 |F(\theta)|\le M_0^{1-\theta}M_1^\theta.
\]
\end{lemat}
\begin{proof}
Najpierw zakładamy $M_0,M_1>0$. Funkcja
$G(z)=F(z)\exp(-(1-z)\log M_0-z\log M_1)$
jest ograniczona na pasie i ma moduł co najwyżej jeden
na jego pionowych brzegach. Dla $\varepsilon>0$ połóżmy
$H_\varepsilon(z)=G(z)\exp(\varepsilon(z^2-1))$.
Jeżeli $z=x+iy$, $0\le x\le1$, to moduł dodatkowego
czynnika wynosi $\exp(\varepsilon(x^2-y^2-1))$.
Na pionowych brzegach jest więc nie większy od jedności.
Na poziomych odcinkach $y=\pm R$ moduł $H_\varepsilon$
nie przekracza $\|G\|_\infty e^{-\varepsilon R^2}$,
co dla dużych $R$ jest mniejsze od jedności.
Z zasady maksimum modułu na prostokącie otrzymujemy
$|H_\varepsilon(\theta)|\le1$.
Przejście do granicy $\varepsilon\downarrow0$ daje
$|G(\theta)|\le1$. Dla $M_j=0$ stosujemy udowodniony
wynik ze stałymi $M_j+\delta$ i przechodzimy do
granicy $\delta\downarrow0$.
\end{proof}

\begin{twierdzenie}[Twierdzenie Riesza--Thorina]
\label{thm:dt-riesz-thorin}
Niech $(X,\mu)$ i $(Y,\nu)$ będą przestrzeniami
z miarami $\sigma$-skończonymi. Niech $T$ będzie
operatorem liniowym określonym na zespolonych funkcjach
prostych o nośniku skończonej miary w $X$, o wartościach
w klasach funkcji mierzalnych w $Y$.
Załóżmy, że dla $j=0,1$ oraz wszystkich takich $f$
\[
 \|Tf\|_{q_j}\le M_j\|f\|_{p_j},
 \qquad 1\le p_j<\infty,\quad 1\le q_j\le\infty.
\]
Dla $0<\theta<1$ definiujemy $p,q$ przez
\[
 \frac1p=\frac{1-\theta}{p_0}+\frac{\theta}{p_1},
 \qquad
 \frac1q=\frac{1-\theta}{q_0}+\frac{\theta}{q_1},
\]
przyjmując $1/\infty=0$. Wówczas $T$ przedłuża się
jednoznacznie do ograniczonego operatora
$L^p(X)\to L^q(Y)$, którego norma nie przekracza
$M_0^{1-\theta}M_1^\theta$.
\end{twierdzenie}
\begin{proof}
Przyjmijmy najpierw dodatnie $M_j$. Wystarczy rozważyć
niezerowe funkcje proste $f,g$ o nośnikach skończonej
miary, unormowane przez $\|f\|_p=\|g\|_{q'}=1$.
Na zbiorach, na których funkcje są niezerowe, piszemy
$f=\omega_f|f|$ oraz $g=\omega_g|g|$, gdzie
$|\omega_f|=|\omega_g|=1$. Dla $z$ w pasie definiujemy
\[
\begin{aligned}
 a(z)&=p\bigl((1-z)/p_0+z/p_1\bigr),\\
 f_z&=\omega_f\exp(a(z)\log|f|).
\end{aligned}
\]
Poza zbiorem $\{f\ne0\}$ przyjmujemy $f_z=0$.
Jeśli $q'<\infty$, definiujemy analogicznie
\[
 b(z)=q'\bigl((1-z)/q_0'+z/q_1'\bigr),
 \qquad
 g_z=\omega_g\exp(b(z)\log|g|),
\]
z wartością zero poza $\{g\ne0\}$.
Jeżeli $q'=\infty$, to $q=q_0=q_1=1$ i bierzemy
po prostu $g_z=g$.
Mamy $f_\theta=f$, $g_\theta=g$. Na brzegach pasa
$\|f_{j+it}\|_{p_j}=1$. Istotnie, podniesienie
modułu do potęgi $p_j$ daje $|f|^p$.
Podobnie $\|g_{j+it}\|_{q_j'}\le1$; jeśli
$q_j'=\infty$ w przypadku $q'<\infty$, moduł
na niezerowym nośniku jest równy jeden.

Rozważmy biliniowe parowanie
\[
 F(z)=\int_Y (Tf_z)(y)g_z(y)\,\mathrm d\nu(y).
\]
Jest ono skończone: rozpisując $f_z,g_z$ na skończone
sumy funkcji charakterystycznych, otrzymujemy całki
$T\chi_A$ po zbiorach skończonej miary, które istnieją
na mocy oszacowań brzegowych i nierówności Höldera.
Liniowość $T$ pokazuje, że $F$ jest skończoną sumą
funkcji wykładniczych o wykładnikach liniowych w $z$.
Ponieważ współczynniki przy $z$ w tych wykładnikach
są rzeczywiste, $F$ jest holomorficzna i ograniczona
na pasie. Na jego brzegach, z nierówności Höldera,
\[
 |F(j+it)|\le
 \|Tf_{j+it}\|_{q_j}\|g_{j+it}\|_{q_j'}\le M_j.
\]
Lemat~\ref{lem:dt-three-lines} daje
$|\int_Y (Tf)g\,\mathrm d\nu|\le M_0^{1-\theta}M_1^\theta$.
Po uwzględnieniu dowolnych norm $f,g$ otrzymujemy
\[
 \left|\int_Y(Tf)g\,\mathrm d\nu\right|
 \le M_0^{1-\theta}M_1^\theta\|f\|_p\|g\|_{q'}.
\]

Wyjaśnijmy, dlaczego wystarcza testowanie funkcjami
prostymi. Niech $u$ będzie funkcją mierzalną skończoną
prawie wszędzie, lokalnie całkowalną na zbiorach
skończonej miary, oraz niech
$|\int ug\,\mathrm d\nu|\le C\|g\|_{q'}$ dla takich
funkcji prostych $g$. Jeżeli $1<q<\infty$, wybieramy
rosnące wyczerpanie $Y$ zbiorami $B_n$ skończonej miary
i kładziemy $E_n=B_n\cap\{|u|\le n\}$. Testujemy
ograniczoną funkcją
$\chi_{E_n}\overline{\omega_u}|u|^{q-1}$, z fazą
$\omega_u=u/|u|$ na $\{u\ne0\}$ i zerem w pozostałych
punktach. Jest ona jednostajną granicą funkcji prostych
o nośnikach w $E_n$, więc nierówność przechodzi do
granicy zarówno w całce, jak i w normie $L^{q'}$.
Dostajemy
\[
 \int_{E_n}|u|^q\,\mathrm d\nu
 \le C\left(\int_{E_n}|u|^q\,\mathrm d\nu\right)^{1/q'}.
\]
Wynika stąd $\|\chi_{E_n}u\|_q\le C$, a twierdzenie
o zbieżności monotonicznej daje $\|u\|_q\le C$.
Dla $q=1$ stosujemy test $\chi_{E_n}\overline{\omega_u}$
o normie supremum nie większej od jedności.
Dla $q=\infty$ i $\eta>0$, gdyby zbiór
$\{|u|>C+\eta\}$ miał dodatnią miarę, miałby podzbiór
$B$ o dodatniej skończonej mierze. Test
$\chi_B\overline{\omega_u}$, znów przybliżony jednostajnie
funkcjami prostymi, prowadziłby do sprzeczności
$(C+\eta)\nu(B)\le C\nu(B)$.
Stosując ten argument do $u=Tf$, otrzymujemy żądane
oszacowanie normy. Przypadek zerowego $M_j$ wynika
z zastąpienia go przez $M_j+\delta$.
Ponieważ $p<\infty$, funkcje proste o nośniku skończonej
miary są gęste w $L^p$, a $L^q$ jest zupełna.
Przedłużenie przez granice jest zatem dobrze określone,
ograniczone i jednoznaczne.
\end{proof}

\begin{twierdzenie}[Nierówność Hausdorffa--Younga]
\label{thm:dt-hausdorff-young}
Niech $1\le p\le2$, a $p'$ będzie wykładnikiem sprzężonym.
Dla unitarnej normalizacji transformacji Fouriera
z Definicji~\ref{def:fourier-schwartz17} zachodzi
\[
 \|\widehat f\|_{p'}
 \le (2\pi)^{-d(1/p-1/2)}\|f\|_p,
 \qquad f\in L^p(\R^d).
\]
Transformata jest jednoznacznym ciągłym przedłużeniem
operatora całkowego z funkcji prostych o nośniku
skończonej miary, zgodnym z transformacją na $\mathcal S'$.
\end{twierdzenie}
\begin{proof}
Dla $p=1$ stosujemy Twierdzenie~\ref{thm:dt-fourier-l1-rl},
dla $p=2$ Twierdzenie~\ref{thm:plancherel22}.
Dla $1<p<2$ korzystamy z Twierdzenia~\ref{thm:dt-riesz-thorin}
z parametrami
\[
 (p_0,q_0,M_0)=(1,\infty,c_d),\qquad
 (p_1,q_1,M_1)=(2,2,1).
\]
Oba oszacowania dotyczą tego samego operatora dzięki
Stwierdzeniu~\ref{prop:dt-fourier-l1-l2-compat}.
Mamy $1/p=1-\theta/2$, $1/q=\theta/2$, zatem
$q=p'$ i $1-\theta=2/p-1$.
Norma interpolowana wynosi
$c_d^{\,2/p-1}=(2\pi)^{-d(1/p-1/2)}$.

Aby sprawdzić zgodność z dystrybucjami, wybieramy
funkcje proste $f_n$ o nośnikach skończonej miary
zbiegające do $f$ w $L^p$. Wówczas
$\widehat f_n\to\widehat f$ w $L^{p'}$.
Dla $\varphi\in\mathcal S$ nierówność Höldera daje
\[
 \left|\int_{\R^d}(f_n-f)\mathcal F\varphi\,
 \mathrm d\lambda_d\right|
 \le\|f_n-f\|_p\|\mathcal F\varphi\|_{p'}\to0,
\]
a analogicznie parowanie $\widehat f_n-\widehat f$
z $\varphi$ dąży do zera. Każde $f_n$ należy do $L^1$,
więc stosujemy
Stwierdzenie~\ref{prop:dt-fourier-l1-l2-compat}
i przechodzimy do granicy w tożsamości biliniowej.
Dostajemy $\mathcal FT_f=T_{\widehat f}$.
\end{proof}

\begin{uwaga}[Zakres i stała]
\label{rem:dt-hy-constant}
Powyższa stała pochodzi z interpolacji dwóch norm
brzegowych. Nie twierdzimy, że jest optymalna dla
$1<p<2$. Ostre oszacowanie Babenki--Becknera
wymaga odrębnego dowodu.
Nie można natomiast rozszerzyć samego oszacowania
$\|\widehat f\|_{p'}\le C\|f\|_p$ na $p>2$
przez formalną zamianę wykładników. Kontrprzykład
do takiego rozszerzenia zostanie omówiony poniżej.
\end{uwaga}

\begin{stwierdzenie}[Brak analogicznego oszacowania dla $p>2$]
\label{prop:dt-hy-fails}
Dla $2<p\le\infty$ nie istnieje stała $C$ taka, że
$\|\widehat f\|_{p'}\le C\|f\|_p$ dla wszystkich
$f\in\mathcal S(\R^d)$.
\end{stwierdzenie}
\begin{proof}
Niech $\psi$ będzie niezerową funkcją gładką
o nośniku w kuli o promieniu $1/4$ i niech
$\phi=\mathcal F^{-1}\psi\in\mathcal S$.
Dla $N\ge1$ i znaków $\varepsilon_j\in\{-1,1\}$
kładziemy
\[
 f_N(x)=\phi(x)\sum_{j=1}^N\varepsilon_j e^{ijx_1}.
\]
Transformaty składników mają parami rozłączne nośniki,
więc dla $q=p'<2$ (także $q=1$)
\[
 \|\widehat f_N\|_q=N^{1/q}\|\psi\|_q.
\]
Dla skończonego $p>2$ wybierzmy znaki niezależnie,
z równymi prawdopodobieństwami obu wartości.
Dla $m$ całkowitego z $2m\ge p$ rozwijamy
$\left|\sum_j\varepsilon_j e^{ijx_1}\right|^{2m}$.
Wartość oczekiwana składnika znika, jeśli którykolwiek
indeks występuje nieparzystą liczbę razy. Pozostałe
składniki mają moduł jeden i najwyżej $m$ różnych
indeksów. Ich liczba nie przekracza
$(2m)^{2m}N^m$ (wybór wzorca indeksów i samych indeksów).
Nierówność Höldera na przestrzeni probabilistycznej daje
zatem
\[
 \mathbb E\left|\sum_j\varepsilon_j e^{ijx_1}\right|^p
 \le (2m)^p N^{p/2}.
\]
Całkując i stosując Tonellego, znajdujemy wybór znaków
taki, że $\|f_N\|_p\le 2m\sqrt N\|\phi\|_p$.
Gdyby rozważane oszacowanie zachodziło, to
$N^{1/q-1/2}$ byłoby ograniczone niezależnie od $N$,
co jest niemożliwe, gdy $q<2$.

Dla $p=\infty$ korzystamy z jawnej konstrukcji
wielomianów Rudina--Shapira:
$P_0=Q_0=1$ oraz
\[
 P_{n+1}(z)=P_n(z)+z^{2^n}Q_n(z),\qquad
 Q_{n+1}(z)=P_n(z)-z^{2^n}Q_n(z).
\]
Indukcyjnie oba mają po $2^n$ współczynników
równych $1$ lub $-1$, ponieważ nośniki potęg
w obu składnikach są rozłączne. Dla $|z|=1$ mamy
$|P_{n+1}|^2+|Q_{n+1}|^2=2(|P_n|^2+|Q_n|^2)$,
więc $|P_n(z)|\le\sqrt{2^{n+1}}$.
Biorąc $N=2^n$ i $f_N(x)=\phi(x)e^{ix_1}P_n(e^{ix_1})$,
otrzymujemy $\|f_N\|_\infty\le\sqrt{2N}\|\phi\|_\infty$,
podczas gdy $\|\widehat f_N\|_1=N\|\psi\|_1$.
To również wyklucza istnienie stałej $C$.
\end{proof}
